\section{Implementación}

TuEvento es una plataforma completa de gestión de eventos que permite a usuarios finales descubrir y asistir a eventos, mientras que organizadores pueden crear y gestionar sus eventos. El sistema incluye funcionalidades de autenticación, perfiles de usuario, búsqueda de eventos, compra de boletos, calificaciones y notificaciones.

\subsection{Arquitectura General}
La arquitectura sigue el patrón MVC estándar, con separación clara entre las capas:
\begin{itemize}
\item \textbf{Model}: Gestiona datos y lógica de negocio (entidades JPA, repositorios, servicios).
\item \textbf{View}: Presenta la interfaz de usuario (páginas web en React, pantallas móviles en React Native).
\item \textbf{Controller}: Maneja solicitudes HTTP y coordina entre Model y View.
\end{itemize}

El sistema consta de tres componentes principales:
\begin{enumerate}
\item \textbf{Backend (Spring Boot)}: API REST que expone endpoints para todas las operaciones.
\item \textbf{Frontend Web (React)}: Interfaz web responsiva para usuarios de escritorio.
\item \textbf{Frontend Móvil (Expo/React Native)}: Aplicación móvil nativa.
\end{enumerate}

\subsection{Backend: Spring Boot}
El backend está desarrollado con Spring Boot 3.3.4, utilizando Java 17. Las dependencias principales incluyen:
\begin{itemize}
\item Spring Data JPA para persistencia.
\item Spring Security con JWT para autenticación.
\item Spring Mail para envío de emails.
\item AWS SDK S3 para almacenamiento de archivos.
\item Apache POI para procesamiento de Excel.
\end{itemize}

\subsubsection{Modelo de Datos}
El modelo de datos incluye las siguientes entidades principales:
\begin{itemize}
\item \textbf{User}: Información de usuarios con roles (ADMIN, USER), direcciones y estado.
\item \textbf{Event}: Eventos con organizador, ubicación, fechas y categorías.
\item \textbf{Ticket}: Boletos asociados a usuarios y eventos.
\item \textbf{Seat/Section}: Asientos organizados en secciones con precios.
\item \textbf{EventRating}: Calificaciones y comentarios de usuarios.
\item \textbf{Notification}: Sistema de notificaciones por evento.
\end{itemize}

\subsubsection{Capa de Controladores}
Los controladores REST implementan endpoints para:
\begin{itemize}
\item Gestión de usuarios (registro, login, perfiles).
\item CRUD de eventos.
\item Compra y gestión de boletos.
\item Calificaciones y reseñas.
\item Administración (gestionar usuarios, peticiones de organizadores).
\end{itemize}

\subsubsection{Capa de Servicios}
Los servicios contienen la lógica de negocio, incluyendo:
\begin{itemize}
\item Validación de datos.
\item Procesamiento de pagos (simulado).
\item Envío de notificaciones.
\item Generación de códigos QR para boletos.
\end{itemize}

\subsection{Frontend Web: React}
El frontend web utiliza React 19 con Vite como bundler. Las tecnologías incluyen:
\begin{itemize}
\item React Router para navegación.
\item Tailwind CSS para estilos.
\item Framer Motion para animaciones.
\item Axios para llamadas API.
\end{itemize}

\subsubsection{Estructura de Componentes}
\begin{itemize}
\item \textbf{Páginas principales}: Landing page, lista de eventos, detalles de evento, perfil de usuario.
\item \textbf{Componentes reutilizables}: Botones, inputs, tarjetas de evento.
\item \textbf{Hooks personalizados}: Para gestión de estado y llamadas API.
\end{itemize}

\subsubsection{Flujo de Usuario}
\begin{enumerate}
\item Usuario se registra/verifica email.
\item Explora eventos por categoría o ubicación.
\item Selecciona evento y compra boletos.
\item Recibe confirmación y código QR.
\item Puede calificar el evento después de asistir.
\end{enumerate}

\subsection{Frontend Móvil: Expo/React Native}
La aplicación móvil utiliza Expo SDK con React Native. Incluye:
\begin{itemize}
\item React Navigation para navegación (stack y tabs).
\item NativeWind para estilos consistentes.
\item Expo Camera para escaneo QR.
\item AsyncStorage para persistencia local.
\end{itemize}

\subsubsection{Pantallas Principales}
\begin{itemize}
\item \textbf{Auth}: Login, registro, verificación.
\item \textbf{Eventos}: Lista, búsqueda, detalles.
\item \textbf{Perfil}: Información personal, tickets.
\item \textbf{QR Scanner}: Para validar boletos.
\end{itemize}

\subsection{Integración y Despliegue}
\begin{itemize}
\item \textbf{Base de datos}: PostgreSQL para datos relacionales.
\item \textbf{Almacenamiento}: AWS S3 para imágenes de eventos.
\item \textbf{Email}: Gmail SMTP para notificaciones.
\item \textbf{Despliegue}: Docker Compose para desarrollo local.
\end{itemize}

El sistema demuestra cómo MVC facilita el desarrollo paralelo de equipos especializados en backend, frontend web y móvil, manteniendo consistencia en la arquitectura y APIs.